\documentclass[../../main2.tex]{subfiles}

\begin{document}

\chapter{目标：把意向写成可比较的量}
\label{ch:goal}

\begin{motivation}
	上一章遗留的问题：动机不决定结果，因为从动机到结果隔着「目标」和「环境」两层转换。环境留给 Part II；本章处理前一层——但「目标」还只是日常词汇：「我想要 X」，X 到底是什么形状？一个名词？一张清单？一个数？\\
	\textbf{核心动机}：把「目标」从愿望改写成可比较、可反推的东西。没有这一步，上一节说的「落差」只是修辞，预测无从谈起。
\end{motivation}

\begin{logicmap}
\begin{tikzpicture}[node distance=0.85cm, every node/.style={draw, rectangle, rounded corners, align=center, font=\small}]
\node (q) {「想要 X」：X 是什么形状？};
\node (rank) [below=of q] {目标 = 排序，不是清单};
\node (comm) [below=of rank] {尺子不止一把：不可通约};
\node (rev) [below=of comm] {从行为反推：显示性偏好};
\node (next) [below=of rev] {反推不唯一，退到概率断言};
\draw[-{Stealth}] (q) -- (rank);
\draw[-{Stealth}] (rank) -- (comm);
\draw[-{Stealth}] (comm) -- (rev);
\draw[-{Stealth}] (rev) -- (next);
\end{tikzpicture}
\end{logicmap}

\section{目标不是名词，是排序}

\begin{readingnote}
	你有没有想过，为什么「我的目标是有钱、有闲、有健康」这种话，说的时候很爽，说完等于没说？\\
	因为这三样\textbf{打架}的时候你听谁的，才是你的目标——清单本身不是。
\end{readingnote}

先做一个残酷但诚实的观察：没有人只「想赚钱」。一个人或者想要「钱多过现在两倍」，或者想要「钱够花但六点下班」，或者想要「钱比同学多」。「想赚钱」不是目标，是\textbf{目标集合}——集合不驱动行为，取舍才驱动行为。

场景永远是取舍：加班还是回家？跳槽还是留下？一个人在冲突时刻的\textbf{顺序}，才是他的目标。

\begin{quote}
	目标的正确形状不是名词，是\textbf{排序}（preference ordering）——在任意两个选项之间，排出先后。
\end{quote}

用本书的标准句式把它分离出来：

\begin{quote}
	如果「想要」不需要比较，那它只是情绪——一阵想要，一阵不想要，预测不了任何行为。但现在加入一个约束：\textbf{目标必须能指导选择}——能指导选择的东西，必须能在任意两个选项之间排出先后。这一步改变了所有性质，这就是「目标」从「愿望」中分离出来的时刻。
\end{quote}

例子就摆在每个人的生活里：

\begin{itemize}
	\item \textbf{买房的人}嘴上说「要近地铁、要好学区、要大平米、要便宜」——全都要。预算一压，中介只带他看两套房：A 近地铁六十平，B 远地铁一百平学区好。他选了哪个，他的目标就写在哪个选择里。嘴上的清单是广告，行为上的排序才是本人。
	\item \textbf{选课的学生}说「都想学」。时间冲突那天，他选了那门课——排序现身了。
	\item \textbf{接 offer 的上班族}说「钱、成长、稳定都重要」。签下去的那家公司，暴露了权重。
	\item \textbf{古代的和亲}：王朝说「和平、面子、土地都重要」。真到取舍那天，送出去的是公主，留下来的边疆——排序写在史书的夹缝里。
\end{itemize}

\begin{questionbox}
	\textit{现在你可能想反驳：「排序太苛刻了吧？中午吃面还是吃饭，我真心无所谓。」}\\
	\textbf{无所谓本身就是排序的一部分}：平局是合法的。真正破坏排序的是\textbf{循环}：面比饭好，饭比面包好，面包又比面好。行为经济学证明循环偏好真实存在（换个说法就能当场把偏好拧反，叫框架效应（framing effect））。\lowfreq{本书的工作假设是：在生存相关的要紧事上，偏好近似可排序}——小事随缘，大事不绕圈。
\end{questionbox}

\paragraph*{逻辑衔接}
	排序解决了「能不能比」。但「比」需要尺子——所有选项都放在同一把尺子上量吗？下一节给坏消息：尺子不止一把，而且有些尺子之间\textbf{没有换算率}。

\section{尺子不止一把}

\begin{readingnote}
	你有没有想过，「这条命值多少钱」这个问题，为什么让人本能地不舒服？\\
	不是问题残忍——是你心里知道：\textbf{有些东西换算不成别的}。
\end{readingnote}

三样东西放一起：金钱、自由、生命。经济学教科书喜欢假装一切皆有价。做个思想实验就能摸到边界：保险公司确实为生命定价，但注意它定的是什么——是\textbf{降低一点死亡风险}的价格，不是\textbf{命}的价格。你愿意花两千块买保险降低一点风险，和你愿意收多少钱\textbf{去死}，是两个宇宙里的问题。越过某条线，「报价」这个动作本身失灵。

本书在这一节的立场，不是折中：

\begin{quote}
	\textbf{目标之间不能完全通约。且这不是人类理性的缺陷——这恰恰是政治与道德存在的理由。}
\end{quote}

反过来想就通了：如果生命、自由、金钱之间真有完备的换算表，一切价值冲突都只是计算错误，人类就不需要道德争论，只需要\textbf{会计}。正因为换算表不存在，排序冲突才无解于算盘，人类才需要超出计算的装置来裁决——伦理是它，政治也是它。

\begin{tcolorbox}[colback=green!5, title=\textit{生活类比}]
	多把尺子的世界像体检报告：血压、血糖、体重各是一把尺。你不能说「血压 120 换算成体重是多少斤」。医生怎么综合？靠经验加权——吵架的时候，就是不同医生把不同尺子排在了前面。
\end{tcolorbox}

跨人群的「不可通约」现场，你肯定见过至少一个：

\begin{itemize}
	\item \textbf{学生}：「离家近的好学校」对「离家远的顶尖专业」——名次和亲情，没有官方汇率。
	\item \textbf{上班族}：「高薪但 996」对「低薪但接孩子」——钱和陪伴，换算率每天都在变。
	\item \textbf{做研究的}：「稳定教职」对「高风险方向」——安稳和好奇，不进同一张表。
	\item \textbf{历史里}：文天祥面前是「活命」对「名节」；丝绸之路上的商队是「利润」对「沙暴里的命」。每个时代都在给不可通约的事物找临时汇率，找不到就动刀兵。
\end{itemize}

\begin{questionbox}
	\textit{现在你可能想反驳：「经济学不就假设人人追求效用最大化吗？那不就是假定一切都通约了？」}\\
	\textbf{把两件事分开：追求最大和单一尺度是两回事。}人完全可以对多把尺子排序一致，而不把它们折成钱。把经济学读成拜金教程，是把工具当成了世界观。
\end{questionbox}

\begin{reader_anticipation}
	\item 不可通约的冲突，最后由谁裁决？→ 第 13 章「政治」：有些偏好没法标价，政治就是接手这个裁决的东西。
	\item 「各占多少」会不会也遇到不可加？→ 会。第 9 章和数学附录：两个原因加起来不等于总效应，交互部分必须单独处理——不可通约性在因果空间里的影子。
\end{reader_anticipation}

\paragraph*{逻辑衔接}
	尺子不止一把，可现实的选择总是「二选一」——多维的想要必须被压回一维才能行动。这个压缩动作，工程师其实天天在做：打分、加权、排序。但权重写在谁的脑子里？下一节：从行为反推。

\section{从行为反推：显示性偏好}

\begin{readingnote}
	你有没有想过，为什么「听其言」之后还要「观其行」？\\
	因为嘴上的目标可以无限美好，行为会把真实的排序\textbf{供}出来。
\end{readingnote}

经济学家萨缪尔森把这个直觉锻造成了工具：\textbf{显示性偏好}（revealed preference）。逻辑一句话：一个人在能选 A 的情况下选了 B，说明对他而言 B 不比 A 差。不用读心，只用记账——账本就是他的选择史。

例子：同事口头宣称「健康最重要」，却连续三年接受每天通勤三小时的高薪岗位。嘴和脚给了两份不同的权重表。显示性偏好只认脚。

\begin{quote}
	如果权重藏在人心里看不见，反推就不可避免——目标函数从「给定」变成「推断」，从描述问题变成统计问题。从此，「想要什么」永远隔着数据的一层纱。
\end{quote}

然后是本章第二句重要的诚实话：\textbf{反推不唯一}。买咖啡这个行为，跟三个互不相同的动机都兼容：真喜欢味道；需要咖啡因提神；办公室文化让不喝咖啡的人显得不合群。三种动机，同一张消费小票。这不是数据不够多，是\textbf{结构上分不开}。

单个行为只能给出排序的\textbf{下界}：排除「被否决的选项排在所选之上」，仅此而已。要收窄，只有一条路：让同一个人在不同约束下反复选择——价格变、环境变、同伴变，每次切掉一块候选。但无论切多少刀，有限次观测永远拟合不出唯一的答案。

\begin{questionbox}
	\textit{现在你可能想反驳：「框架效应能把偏好当场拧反——反推出来的排序还作数吗？」}\\
	\textbf{作数，但换个读法：局部排序。}反推出的排序只在观测到的约束范围内有效，出了这个范围不作数。偏好随状态变——这不是让步，是预告：状态会塑造偏好，这件事在第 16 章「循环的闭合」里还会回来咬你一口。
\end{questionbox}

\paragraph*{逻辑衔接}
	清点家当：目标是排序；尺子不止一把；从行为能反推但不唯一。加起来得到一个看似让步、实为全书转折点的结论——下一节收束。

\section{交给概率：预测的正确形状}

\begin{readingnote}
	你有没有想过，为什么天气预报不说「明天降水」，只说「降水概率七成」？\\
	不是气象学家怯懦——是他们对「预测」这件事\textbf{更诚实}。
\end{readingnote}

把本章四步连成一句话：关于「人想要什么」，我们永远拿不到确定的答案——能拿到的，只有「在给定条件下，行为与结果以多大概率发生」。

这个结论规定了预测的数学形状：\textbf{预测是概率断言，不是预言}。说「他有七成概率跳槽」永远不会被单个案例证伪——可检验性藏在多次、多对象的统计里。这就是导论那句「气象员与巫师干的是两种活」的技术版本。

但概率不能凭空落地。它需要一个落脚点：「未来由什么构成」必须先被说清，概率才有地方站——往什么东西赋概率？这正是 Part II 的任务。而且这条依赖是真的：

\begin{quote}
	概率必须定义在「世界有哪些状态」上。不先说清未来由什么构成，概率就是空转的符号。
\end{quote}

\begin{questionbox}
	\textit{现在你可能想反驳：「绕了一大圈，结论就是『算概率』？开头直接讲概率不就完了？」}\\
	\textbf{不行——不经过目标这一站，你不知道概率该定义在什么上。}「对什么的概率？在哪些条件下？」答案埋在本章：目标是排序，行为是它的投影，所以可观测的是\textbf{行为的概率}。跳过这一站直接写概率的书，会在「概率定义在什么上」摔跟头。
\end{questionbox}

\paragraph*{逻辑衔接}
	本章把「想要」从名词改成了排序，又诚实地退到了概率。Part I 到此闭合，但有一个大问题必须在进入 Part II 之前，当着你的面埋下——

\section{Part I 尾声：一个埋给 Part IV 的追问}

一个人的目标清楚了，行为也大致能反推了。可是社会不是一个人。千万人的排序搅在一起，会自动变成稳定的社会规律吗——菜价、房价、舆论走向，是千万个动机的\textbf{加总}，还是有它自己的脾气？

\begin{warningbox}
	这是本书第一枚定时的追问：\textbf{微观的动机，怎么汇聚成宏观的结构？}\\
	现在不答。Part II 先把「预测与因果」的工具备齐，Part IV 开头（第 10 章「通道生成器」）正面回答它。成对出现，不许只埋不收——这是全书的规矩。
\end{warningbox}

\paragraph*{逻辑衔接}
	工具备到这一步：人有动机，动机化为排序，排序只能从行为反推，反推的结果只能是概率。下一章问那个导论钉下的枢纽问题：为什么「能预测」和「有因果」，是同一枚硬币的两面？

\begin{thinkbox}
\begin{enumerate}
	\item 翻出你最近三个月的三笔「非必要」消费（一杯奶茶、一次打车、一个会员）。每笔背后是哪个排序赢了？把它们连起来，画出你嘴上没说过的权重表。
	\item \textbf{反事实}：如果「钱、自由、生命」之间真的存在官方换算表——一切皆可标价——世界会变成什么样？法庭还存在吗？政治还存在吗？用本章「不可通约是政治存在的理由」推一遍。
	\item \textbf{反事实}：如果你的公司在你入职那天公开了全部晋升打分公式，你的行为会怎么变？「从行为反推目标」这门手艺，会变容易还是变难？
\end{enumerate}
\end{thinkbox}

\end{document}
